\section{Motivation and Objectives}
\todo{Forse dovrei concentrarmi più su RNN dall'inizio?}
Nowadays, Artificial Intelligence (AI) became the main driver of innovation in many fields, at the point that the growth of research
became exponential in the last decade. This fully alignes with the current state of our world, in fact,
we are observing a paradigm shift in the way we interact with technology (or world), through the use of AI.\newline
Within AI, whose main goal is not only to understand \textit{intelligence}, but to embed it into agents capable of performing complex tasks through
\textit{rational reasoning, planning} and \textit{acting}, the subfield of Machine Learning (ML) has emerged as the dominant approach to mimicking biological intelligent behaviour.
Neural Networks (NN) \textit{learn} complex patterns from large amounts of real-world data, instead of relying on explicit rule based approaches,
in order to perform \textit{inference} and return prediction on unseen data. \newline
The last years have seen a surge in popularity of generative models, after a first wave of discriminative models that made popular the field of Deep Learning (DL),
and introduced GPU computation in 2009 \cite{Raina2009LSDL}. 
Particularly, deep generative models have been applied with unbelievable results in many fields,
 from image generation with \textit{Diffusion Models} \cite{ho2020denoisingdiffusionprobabilisticmodels}
 to text generation with \textit{Transformers} \cite{vaswani2023attentionneed} and protein folding with
\textit{Alphafold} \cite{jumper2021highly} and many others domains, achieving state of the art results.
However, deep learning models are known to be expensive in both training and inference, requiring large 
amount of computational resources (GPU production) and energy consumption \cite{article}.

\section{Reservoir Computing for Efficient AI} \todo{mancano citazioni ed è ripetitivo rispetto al background}
Sequential data processing remains a fundamental challenge in machine learning, from natural language to time series forecasting.
While Recurrent Neural Networks (RNNs) have been the traditional approach for modeling temporal dependencies, 
training them to capture long-term patterns has proven notoriously difficult due to vanishing and exploding gradients.
Although architectures like LSTMs and GRUs have partially addressed these issues through gating mechanisms,
and Transformers have recently dominated the field with their attention-based approach, these solutions come at the cost
of increased architectural complexity and computational demands.\newline
An alternative paradigm that has gained attention is Reservoir Computing (RC)\cite{}, which offers a fundamentally different approach
to sequential processing. By maintaining a fixed, randomly initialized recurrent network as a "reservoir" of dynamics,
and training only a simple linear readout layer, RC circumvents the gradient-based training of recurrent connections entirely.
This framework not only simplifies the learning process but also draws inspiration from neuroscience and physical computing systems,
where rich dynamics emerge naturally without explicit optimization.\newline

Despite its computational efficiency and theoretical elegance, vanilla RC models face their own set of limitations,
particularly in scaling to complex tasks that require multiple levels of temporal abstraction\cite{}.
This motivates the exploration of \textit{deep} variants \cite{} of reservoir computing architectures,
which aim to combine the efficiency of the RC paradigm with the representational power of hierarchical processing.
This thesis investigates such deep reservoir architectures, examining how stacking multiple reservoir layers
can enable learning at different timescales while maintaining the training efficiency that makes RC attractive.\newline


\section{Scientific contribuition}
\todo{After finishing experimental part...}


\section{Structure of the work}

elenco puntato con quello che si dirà in ogni capitolo della tesi